\documentclass[11pt,a4paper]{article}
\usepackage[utf8]{inputenc}
\usepackage[T1]{fontenc}
\usepackage[portuguese]{babel}
\usepackage{xcolor}
\usepackage{hyperref}
\usepackage{geometry}
\usepackage{listings}
\usepackage{tcolorbox}
\geometry{margin=2.5cm}
\definecolor{primary}{RGB}{41,128,185}
\definecolor{secondary}{RGB}{44,62,80}
\definecolor{codebg}{RGB}{240,240,240}
\lstset{backgroundcolor=\color{codebg},basicstyle=\ttfamily\small,breaklines=true,frame=single}
\title{\color{secondary}\Huge\textbf{Exercício 3.4: Canary Deployment}}
\author{\color{primary}DevOps com Kubernetes -- 2026}
\date{}
\begin{document}
\maketitle


\textbf{Instruções do Exercício:}

\begin{tcolorbox}[colback=blue!5!white,colframe=blue!75!black]
Planeje uma publicação canário usando labels, annotations ou Argo Rollouts, explicando como validar se a nova versão pode receber mais tráfego.
\end{tcolorbox}

\end{document}
